% =============================================================================
% Capítulo 1 — Introdução
% =============================================================================
\chapter{Introdução}
\label{cap:introducao}

\section{Contextualização}
\label{sec:contextualizacao}

A indústria de jogos digitais consolidou-se como uma das mais relevantes no
cenário do entretenimento contemporâneo, movimentando bilhões de dólares
anualmente e alcançando audiências que superam outras mídias tradicionais
\cite{newzoo2024}. Parte significativa dessa evolução deve-se à sofisticação
crescente dos sistemas de Inteligência Artificial (IA) empregados no controle
de Personagens Não-Jogáveis (NPCs), cuja qualidade comportamental exerce
influência direta sobre a experiência do jogador \cite{millington2019ai}.

Historicamente, a IA em jogos comerciais apoiou-se em técnicas determinísticas,
como Máquinas de Estados Finitos (FSM) e Árvores de Comportamento
(\textit{Behavior Trees}). Essas abordagens oferecem previsibilidade e
facilidade de depuração, qualidades fundamentais em contextos de produção.
Contudo, comportamentos baseados em regras fixas tendem a se tornar
previsíveis após poucas interações, reduzindo o desafio percebido e
comprometendo a rejogabilidade \cite{yannakakis2018artificial}.

O Aprendizado por Reforço (\textit{Reinforcement Learning} --- RL) surgiu como
alternativa promissora para a criação de agentes adaptativos, capazes de
ajustar suas estratégias por meio de interações com o ambiente
\cite{sutton2018reinforcement}. Plataformas como o Unity ML-Agents
\cite{juliani2020unity} democratizaram o acesso a algoritmos de RL no contexto
de motores de jogo, permitindo que desenvolvedores treinem NPCs diretamente
em ambientes simulados. Em particular, o algoritmo \textit{Proximal Policy
Optimization} (PPO) \cite{schulman2017ppo} destacou-se por sua estabilidade
durante o treinamento e por seu bom desempenho em espaços de ação discretos,
características que o tornam adequado para o controle de personagens em jogos
2D.

Entretanto, a adoção de RL em jogos não está isenta de desafios. Agentes
treinados exclusivamente por tentativa e erro podem desenvolver
comportamentos inesperados, difíceis de interpretar e de ajustar por
projetistas de jogo. Nesse contexto, a busca por mecanismos que tornem o
processo de aprendizado mais interpretável --- isto é, que permitam
compreender \textit{por que} o agente aprendeu determinado comportamento ---
constitui uma questão aberta e relevante para a área.

\section{Problema de Pesquisa}
\label{sec:problema}

A definição de funções de recompensa constitui um dos aspectos mais críticos
do treinamento por RL. Quando projetadas manualmente, essas funções tendem a
ser exaustivas, frágeis e difíceis de generalizar: pequenas alterações no
\textit{design} do jogo podem exigir reajustes significativos no esquema de
recompensas \cite{ng1999reward}. Em cenários complexos, a engenharia manual
de recompensas torna-se um gargalo de desenvolvimento.

Uma alternativa amplamente explorada na literatura é a aprendizagem por
imitação. Técnicas como \textit{Generative Adversarial Imitation Learning}
(GAIL) permitem que um agente aprenda a replicar o comportamento de um
jogador humano a partir de demonstrações gravadas. Embora eficaz em diversos
cenários, essa abordagem apresenta uma limitação conceitual relevante: o
agente tende a \textit{copiar} o comportamento observado, sem necessariamente
compreender a relação causal entre as ações executadas e o sucesso obtido. O
resultado é um modelo que reproduz padrões comportamentais sem internalizar a
estratégia subjacente \cite{ho2016gail}.

Emerge, portanto, uma questão de pesquisa: \textbf{é possível guiar o
aprendizado por reforço de um agente de jogo utilizando informações causais
extraídas de sessões de jogo anteriores, sem recorrer à imitação direta de
demonstrações?}

A Proveniência de Dados (\textit{Data Provenance}) oferece um caminho para
endereçar essa questão. Originalmente desenvolvida para rastreamento de
linhagem em bancos de dados e fluxos científicos, a proveniência permite
registrar cadeias causais de eventos --- não apenas \textit{o que} aconteceu,
mas \textit{por que} aconteceu, preservando relações de causa e efeito entre
ações e consequências \cite{provenance_base}. Aplicada ao contexto de jogos
digitais, a proveniência possibilita a construção de grafos que representam a
sequência causal de eventos em uma sessão de jogo, desde ações do jogador até
desfechos como vitória ou derrota \cite{kohwalter2020provenance}.

\section{A Proposta}
\label{sec:proposta}

O presente trabalho propõe uma abordagem que integra Proveniência de Dados e
Aprendizado por Reforço para o treinamento de um agente de jogo. A ideia
central é utilizar grafos de proveniência, gerados a partir de sessões jogadas
contra um chefe controlado por FSM, como fonte de informação causal para guiar
o aprendizado de um agente PPO.

O pipeline proposto opera da seguinte forma:

\begin{enumerate}
    \item Um jogador humano enfrenta o chefe baseline (FSM) em múltiplas
          sessões de jogo.
    \item O sistema de proveniência registra cada evento relevante da sessão
          --- ataques, danos, pulos, esquivas, mortes --- como nós de um
          grafo causal, preservando as relações de causa e efeito entre eles.
    \item Um processo de filtragem seleciona apenas as sessões encerradas
          com vitória do jogador e extrai, a partir dos grafos correspondentes,
          sequências de ações que precederam eventos de dano no chefe.
    \item Essas sequências vitoriosas são convertidas em um artefato
          estruturado e carregadas por um módulo de \textit{Reward Shaping}
          \cite{ng1999reward}, que concede recompensa adicional ao agente PPO
          sempre que este executa padrões de ação compatíveis com os padrões
          extraídos.
    \item O agente é treinado com PPO padrão, utilizando o
          \textit{Reward Shaping} por proveniência como sinal auxiliar. O
          algoritmo continua explorando o ambiente livremente, mas recebe
          orientação indireta das ações humanas bem-sucedidas.
\end{enumerate}

Essa abordagem distingue-se da imitação direta em um aspecto fundamental: o
agente não reproduz trajetórias gravadas nem utiliza um discriminador
adversário para aprender políticas. Em vez disso, ele recebe um sinal de
recompensa adicional que reflete padrões causais associados ao sucesso,
preservando a capacidade exploratória do PPO e mantendo a rastreabilidade das
decisões de \textit{design} do treinamento.

A contribuição acadêmica central, portanto, reside na demonstração de que
sequências extraídas de grafos de proveniência podem ser utilizadas como sinal
auxiliar de recompensa em um agente PPO, produzindo um chefe com comportamento
mensuravelmente distinto do \textit{baseline} FSM em um ambiente de jogo 2D.

\section{Objetivos}
\label{sec:objetivos}

\subsection{Objetivo Geral}
\label{subsec:objetivo-geral}

Desenvolver e avaliar um pipeline de integração entre Proveniência de Dados e
Aprendizado por Reforço (PPO) para o treinamento de um agente de chefe em um
\textit{vertical slice} de jogo 2D, utilizando sequências de ações vitoriosas
extraídas de grafos de proveniência como \textit{Reward Shaping}, e comparando
o comportamento resultante com um chefe \textit{baseline} controlado por
Máquina de Estados Finitos.

\subsection{Objetivos Específicos}
\label{subsec:objetivos-especificos}

Para alcançar o objetivo geral, foram definidos os seguintes objetivos
específicos:

\begin{enumerate}
    \item Implementar o \textit{vertical slice} do jogo \textit{Chronicles Of
          The Lost Word} em Unity, contendo uma arena de chefe com mecânicas
          jogáveis de combate 2D.

    \item Desenvolver um chefe \textit{baseline} controlado por Máquina de
          Estados Finitos (FSM), com comportamento determinístico e
          previsível, adequado para servir como referência de comparação.

    \item Projetar e implementar um sistema de proveniência capaz de registrar
          eventos de jogo como nós de um grafo causal, preservando relações
          de causa e efeito entre ações do jogador e consequências no
          ambiente.

    \item Construir um processo automatizado de extração de sequências de
          ações vitoriosas a partir dos grafos de proveniência gerados em
          sessões contra o chefe FSM.

    \item Implementar um módulo de \textit{Reward Shaping} que carregue as
          sequências extraídas e conceda recompensa adicional a um agente PPO
          durante o treinamento, quando este executar padrões de ação
          compatíveis.

    \item Treinar o agente PPO utilizando o Unity ML-Agents com o
          \textit{Reward Shaping} por proveniência e integrar o modelo
          resultante ao jogo em modo de inferência.

    \item Realizar uma avaliação quantitativa comparativa entre o chefe FSM e
          o chefe PPO, coletando métricas objetivas de desempenho em sessões
          controladas.
\end{enumerate}

\section{Justificativa}
\label{sec:justificativa}

A criação de NPCs com comportamentos interessantes, desafiadores e variados é
um problema recorrente no desenvolvimento de jogos. Embora técnicas
determinísticas como FSM continuem dominantes na indústria, a demanda por
experiências adaptativas motiva a investigação de abordagens baseadas em
aprendizado de máquina \cite{yannakakis2018artificial}.

O Aprendizado por Reforço oferece um arcabouço teórico robusto para esse fim,
mas sua aplicação prática em jogos esbarra em dois obstáculos frequentes: a
dificuldade de projetar funções de recompensa adequadas e a opacidade dos
comportamentos aprendidos. Este trabalho contribui para o enfrentamento de
ambos os obstáculos ao propor o uso de proveniência como ponte entre sessões
de jogo humanas e o treinamento de agentes, introduzindo um mecanismo que é,
ao mesmo tempo, automatizável --- pois não exige ajuste manual de recompensas
para cada cenário --- e interpretável --- pois as sequências que geram
recompensa adicional são derivadas de cadeias causais rastreáveis.

Além disso, a abordagem proposta preenche uma lacuna na literatura ao combinar
duas áreas tipicamente estudadas de forma independente: proveniência de dados
em jogos e aprendizado por reforço para NPCs. A integração entre essas áreas
permite transformar dados de \textit{gameplay} registrados com semântica
causal em sinais de treinamento, sem recorrer a demonstrações diretas ou
técnicas adversárias.

\section{Organização do Trabalho}
\label{sec:organizacao}

O restante deste trabalho está organizado da seguinte forma. O
\autoref{cap:referencial-teorico} apresenta o referencial teórico, abordando
os fundamentos de Máquinas de Estados Finitos, Aprendizado por Reforço, PPO,
Unity ML-Agents, Proveniência de Dados e \textit{Reward Shaping}. O
\autoref{cap:metodologia} descreve a metodologia adotada, incluindo o
pipeline de desenvolvimento, o protocolo de coleta de dados e as decisões de
escopo. O \autoref{cap:implementacao} detalha a implementação técnica do
\textit{vertical slice}, dos sistemas de proveniência e métricas, e da
integração do modelo treinado. O \autoref{cap:avaliacao} apresenta o cenário e
o protocolo de avaliação quantitativa. O \autoref{cap:resultados} expõe os
resultados obtidos na comparação entre os chefes FSM e PPO. O
\autoref{cap:discussao} discute os resultados, analisa as limitações do
trabalho e sugere direções para trabalhos futuros. Por fim, o
\autoref{cap:conclusao} sintetiza as contribuições do trabalho e apresenta as
considerações finais.
